\documentclass[12pt, twoside]{article}
%\documentclass[12pt, twoside]{article}
\usepackage[letterpaper, margin=1in, headsep=0.2in]{geometry}
\setlength{\headheight}{0.6in}
%\usepackage[english]{babel}
\usepackage[utf8]{inputenc}
\usepackage{microtype}
\usepackage{amsmath}
\usepackage{amssymb}
%\usepackage{amsfonts}
\usepackage[nomessages]{fp} %\FPeval{\var-name}{2*sin(pi/6)}
\usepackage{siunitx} %units in math. eg 20\milli\meter
\usepackage{yhmath} % for arcs, overparenth command
\usepackage{tikz} %graphics
\usetikzlibrary{quotes, angles, arrows, arrows.meta}
\usepackage{pgfplots}
\usepgfplotslibrary{fillbetween}
\pgfplotsset{compat=1.16}
\usepackage{graphicx} %consider setting \graphicspath{{images/}}
\usepackage{parskip} %no paragraph indent
\usepackage{enumitem}
\usepackage{multicol}
\usepackage{venndiagram}

\usepackage{fancyhdr}
\pagestyle{fancy}
\fancyhf{}
\renewcommand{\headrulewidth}{0pt} % disable the underline of the header
\raggedbottom
\hfuzz=2mm %suppresses overfull box warnings

\usepackage{hyperref}

\fancyhead[LO]{La Scuola d'Italia / Dr.\ Huson / IB Math 11 \\* 29 September 2026}
\fancyhead[RO]{\fbox{\begin{minipage}{6cm}\footnotesize Nickname:\\[0.5cm]\end{minipage}}}
\fancyhead[LE]{\fbox{\begin{minipage}{6cm}\footnotesize Nickname:\\[0.5cm]\end{minipage}}}
\fancyfoot[C]{\thepage}

\begin{document}
\subsubsection*{1.6 Classwork: Geometric Growth and Decay}

\begin{enumerate}[itemsep=0.15cm]

%--- Problem 1: Common ratio, u_n, u_10, first term over 1 000 000 [6 marks] ---
%--- Source: Original (freshly authored, AA SL 1.3) ---
\item[1.]

In a geometric sequence each term is the previous term multiplied by
the common ratio $r$. The $n^{\text{th}}$ term is
$$u_{n} = u_{1}\,r^{\,n - 1}.$$

Consider the geometric sequence
$$4, \ 12, \ 36, \ 108, \ \ldots$$

\begin{enumerate}[itemsep=0.1cm]
  \item Find $\dfrac{u_{2}}{u_{1}}$, $\dfrac{u_{3}}{u_{2}}$ and
  $\dfrac{u_{4}}{u_{3}}$. Hence write down the common ratio $r$.
  \item Write down an expression for $u_{n}$ in terms of $n$.
  \item Find $u_{10}$.
  \item Find the smallest value of $n$ for which $u_{n} > 1\,000\,000$.
\end{enumerate}

\begin{tikzpicture}
  \draw (-0.6,0) rectangle (14.6,13.4);
  \foreach \y in {1,2,...,13} \draw[dotted] (0,\y)--(14.1,\y);
\end{tikzpicture}

\newpage

%--- Problem 2: Population growth 4% and car depreciation 15% [7 marks] ---
%--- Source: Original (freshly authored, AA SL 1.3, 1.4) ---
\item[2.]

The population of a town was $12\,000$ at the start of 2020. The
population increases by $4\%$ each year.

\begin{enumerate}[itemsep=0.1cm]
  \item Write down the common ratio of the geometric sequence formed by
  the population at the start of each year.
  \item Find the population at the start of 2030.
\end{enumerate}

A car was bought for \$$24\,000$ at the start of 2020. Its value
decreases by $15\%$ each year.

\begin{enumerate}[itemsep=0.1cm]
  \item[(c)] Find the value of the car after $5$ years.
  \item[(d)] Find the number of complete years after which the value
  of the car first falls below \$$5000$.
\end{enumerate}

\begin{tikzpicture}
  \draw (-0.6,0) rectangle (14.6,16.0);
  \foreach \y in {1,2,...,15} \draw[dotted] (0,\y)--(14.1,\y);
\end{tikzpicture}

\end{enumerate}
\end{document}
